\documentclass[12pt]{article}
\usepackage[margin=1in]{geometry}
\usepackage{setspace}
\onehalfspacing
\usepackage[T1]{fontenc}

\title{Regulatory Uncertainty, Developer Entry, and Housing Supply: Evidence from Chicago's Aldermanic Privilege}
\author{Jacob Herbstman}
\begin{document}

\maketitle

At the January 2026 BFI Conference on Budget Analysis and Public Policy, Jeffrey Kling and Heidi Williams delivered a public lecture on how economists can help inform the economic and budget analysis used by Congress. A major theme was the lack of empirical evidence on how permitting and regulatory processes, in particular regulatory uncertainty, affect private investment and development. 

In particular, Heidi Williams highlighted the need for more work on how regulatory uncertainty impacts ``time to build'' and how delays from uncertainty affect the volume and location of private investment.
In December 2025, the CBO published a blog post titled ``A Call for New Research in the Area of Permitting Requirements for Investments in Physical Infrastructure,'' noting that CBO would benefit from additional data and research that would broaden and deepen its basis for modeling the effects of permitting-related legislative proposals. 
In February 2026, the CBO followed up with a report on estimating the effects of changes in permitting requirements on private investment, describing its analytical framework for translating permitting frictions into changes in the capital stock and economic output.

Both documents, along with Williams' speech, highlight that the existing research literature provides very little guidance on some of the key parameters CBO needs, including how permitting delays and uncertainty affect the volume and location of private investment, and how those effects differ across regulatory environments.

This project aims to contribute directly to filling that gap. I will study how aldermanic privilege, the informal norm by which Chicago's City Council defers to the local alderman on land use decisions within their ward, affects housing development. The alderman's support is effectively required at multiple stages of the development approval process, and aldermen vary substantially in how much friction they impose on proposed projects. I aim to build a structural model of this process in which developers choose where to propose projects, anticipating the regulatory cost they will face, and the observed spatial distribution of development reflects the equilibrium of developer location decisions and aldermanic gatekeeping. The data I am seeking to acquire would make this model estimable.

The key empirical challenge is that publicly available data only captures projects that survived the full pipeline. The Chicago Data Portal's building permit dataset, which I currently use, includes only permits that remain valid and explicitly excludes those that were voided, revoked, or withdrawn. This means I can study how aldermanic friction affects the density and scale of completed construction, but I cannot observe the extensive margin: projects that were proposed but killed before reaching the permit stage, or permit applications that were filed but never approved. This is a first-order problem, both for measuring the total effect of regulatory discretion on housing supply and for estimating the structural model, which requires data on the developer's full choice set including locations where projects were attempted but failed.

I am requesting internal administrative records from two City of Chicago departments. From the Department of Buildings, I need the full universe of building permit applications, including those that were denied, withdrawn, voided, revoked, or allowed to expire. The public dataset begins in 2006 and contains roughly 800,000 records, but it reflects only currently valid permits. The complete application file would allow me to measure permit denial and withdrawal rates by ward and year, and to examine whether these rates vary with the alderman in office. Even within the public data, project-level descriptions of proposed work are embedded in unstructured text fields that require scraping and natural language processing to extract usable information about project type, scale, and intended use, so substantial data cleaning work is needed regardless.

From the Department of Planning and Development, I am requesting several categories of records. First, pre-application and concept meeting records. DPD staff conduct concept and land use review meetings with prospective developers before formal applications are filed. If DPD maintains records of these meetings, they would provide a direct measure of the extensive margin of development proposals, which is the most valuable object for the structural entry model. 

Second, the internal zoning application tracking database, which should include all zoning change applications from the point of filing through final disposition, including those withdrawn before the alderman introduced the ordinance to City Council. I have already scraped PDF documents from The City Clerk's Legistar system that documents all rezoning ordinances from 1981 to present day, but even this extensive documentation only captures projects that reached the stage of formal introduction. The gap between DPD's application records and Legistar's legislative records would directly measure how much extensive margin development never comes to fruition. Finally, zoning plan examination records from DPD's Bureau of Zoning, which reviews building permit applications for zoning code compliance before the Department of Buildings conducts its building code review.

Cleaning and processing this data would be an intense and lengthy process. But it would allow me to combine reduced form empirics with structural Industrial Organization (IO) methods to give us the most complete picture yet of how regulatory uncertainty affects investment, specifically in large multifamily and mixed-use developments that are critical to citywide affordability. 

Not only would I be able to extend my current reduced form evidence on how aldermanic privilege affects the density and scale of completed construction (the intensive margin), but I would now be able to measure how it affects the spatial distribution of proposed development and broader economic activity in the city. I would also be able to estimate deeper structural primitives that would allow me to run counterfactual simulations about large policy changes, which would directly inform CBO's modeling of legislative proposals that would change permitting requirements.

I plan to acquire the data through a combination of direct outreach to department leadership, framing the request as an academic research partnership, and FOIA requests where necessary. I have identified the relevant contacts at both departments, including the commissioners and deputy commissioners responsible for the relevant data systems. Some of the records, particularly from DPD, may exist only as unstructured PDFs or scanned documents rather than structured databases, which would require substantial processing and cleaning effort. 

The requested funding would primarily cover the compute costs of extensive OCR processing for thousands of potentially messy PDF documents and scans using the Google Cloud AI API service. I will be processing approximately 200,000 pages of archived city council meetings, along with any other documents I receive from DPD and DOB that are not in structured tabular format.
I will also need funding for cloud storage and compute time for geocoding and spatial analysis of the full dataset, which will likely exceed one million individual records once all applications, meetings, and zoning reviews are combined. 
The timeline will depend heavily on the pace of data acquisition, but I expect to complete the analysis within 12 months of receiving the data.

In addition, I have a separate draft of this paper with only empirical work using publicly-available data, so I will be able to add findings from the new data to that paper at a minimum within that timeframe. 

\end{document}